% ===========================================================================
% Category: Semi-Supervised Learning
% Generated from ML_Lab_Manual_Completed.md - edit the Markdown and re-run
% scripts/build_latex_manual.py instead of editing this file by hand.
% ===========================================================================

\chapter{Experiment 7: Self-training for Semi-supervised
Learning}\label{experiment-7-self-training-for-semi-supervised-learning}

\textbf{Category:} Semi-supervised learning

\section{1. Objective}\label{objective}

Train a classifier with a small labeled subset and a larger unlabeled
subset by iteratively adding high-confidence pseudo-labels; compare with
a supervised baseline and the fully supervised ceiling.

\section{2. Dataset Used}\label{dataset-used}

Breast Cancer Wisconsin (569 × 30, binary). Stratified 70/30 split: 398
training / 171 test. Only 15\% of training labels (61 samples) are kept;
337 samples are masked as unlabeled (−1).

\section{3. Required Libraries}\label{required-libraries}

scikit-learn, numpy

\section{4. Brief Theory}\label{brief-theory}

Self-training fits a base classifier on labeled data, predicts class
probabilities for the unlabeled pool, converts sufficiently confident
predictions (max probability ≥ threshold) to pseudo-labels, adds them to
the labeled set and repeats. Errors can propagate (confirmation bias),
so the confidence threshold and convergence monitoring matter.

\section{5. Procedure}\label{procedure}

\begin{enumerate}
\def\labelenumi{\arabic{enumi}.}
\tightlist
\item
  Split the data stratified; scale features (fit on training data only).
\item
  Mask 85\% of training labels to −1.
\item
  Train a supervised baseline SVC on the labeled 61 samples.
\item
  Wrap an RBF SVC in
  \passthrough{\lstinline!SelfTrainingClassifier(threshold=0.80)!}, fit
  on labeled + unlabeled, monitor iterations.
\item
  Train a fully supervised SVC on all 398 labels as the ceiling.
\item
  Compare accuracy/F1 on the untouched test set.
\end{enumerate}

\section{6. Reference Implementation}\label{reference-implementation}

\begin{lstlisting}[language=Python]
self_training_svc = SelfTrainingClassifier(
    estimator=SVC(kernel='rbf', probability=True, C=1.0, random_state=42),
    threshold=0.80, criterion='threshold', max_iter=15, verbose=True)
self_training_svc.fit(X_train_full_scaled, y_train_semi)   # labels + unlabeled (-1)
\end{lstlisting}

\section{7. Results}\label{results}

{\def\LTcaptype{none} % do not increment counter
\begin{longtable}[]{@{}llll@{}}
\toprule\noalign{}
Setting & Accuracy & F1 & Training samples \\
\midrule\noalign{}
\endhead
\bottomrule\noalign{}
\endlastfoot
Supervised baseline (15\% labeled) & 0.9298 & 0.9469 & 61 \\
Self-training (semi-supervised) & \textbf{0.9415} & \textbf{0.9554} &
382 \\
Fully supervised ceiling (100\%) & 0.9766 & 0.9813 & 398 \\
\end{longtable}
}

Pseudo-labels added per iteration: +292, +22, +5, +2, then termination
\passthrough{\lstinline!no\_change!} (321 new labels total). The
decaying additions indicate stable convergence rather than runaway error
accumulation.

\section{8. Discussion and Conclusion}\label{discussion-and-conclusion}

Using no extra human labels, self-training improves accuracy by +1.2
points over the baseline, closing part of the gap to the fully
supervised model. The remaining gap is the cost of noisy pseudo-labels.
Convergence was quick (4 useful iterations) because the threshold 0.80
admitted mostly correct predictions.

\section{9. Answers to Viva Questions}\label{answers-to-viva-questions}

\begin{itemize}
\tightlist
\item
  \textbf{What is a pseudo-label?} A predicted label for an unlabeled
  sample that the model accepts as training target because its
  confidence exceeds the threshold.
\item
  \textbf{What is confirmation bias in self-training?} The model
  reinforces its own mistakes: wrong confident predictions enter the
  training set, make the model more confident about similar errors, and
  can degrade performance. High thresholds and validation monitoring
  mitigate this.
\item
  \textbf{Why must the test labels remain hidden during training?}
  Otherwise the evaluation is contaminated --- test data would influence
  the model, and the reported accuracy would no longer estimate
  generalization.
\end{itemize}

\begin{center}\rule{0.5\linewidth}{0.5pt}\end{center}

\section{10. Complete Code Listing}\label{complete-code-listing}

\emph{Full runnable program for this experiment, extracted from
\passthrough{\lstinline!notebooks/03\_semi\_supervised\_learning.ipynb!}
(executed; results above are its actual output).}

\textbf{Listing 7.1 --- Suppress a known sklearn deprecation notice for
SVC(probability=True)}

\begin{lstlisting}[language=Python]
# ---- Suppress a known sklearn deprecation notice for SVC(probability=True) ----
import warnings
# ---- Core numerical and plotting libraries ----
import numpy as np
import pandas as pd
import matplotlib.pyplot as plt
import seaborn as sns

# ---- Dataset + train/test utilities ----
from sklearn.datasets import load_breast_cancer
from sklearn.model_selection import train_test_split
from sklearn.preprocessing import StandardScaler
# ---- Semi-supervised self-training wrapper ----
from sklearn.semi_supervised import SelfTrainingClassifier
# ---- Base estimators compared inside self-training ----
from sklearn.svm import SVC
from sklearn.ensemble import RandomForestClassifier
# ---- Evaluation metrics ----
from sklearn.metrics import accuracy_score, precision_score, recall_score, f1_score, roc_auc_score, confusion_matrix, classification_report

# Aesthetics setup (consistent look across all notebooks)
plt.style.use('seaborn-v0_8-whitegrid' if 'seaborn-v0_8-whitegrid' in plt.style.available else 'default')
plt.rcParams['font.sans-serif'] = 'DejaVu Sans'
plt.rcParams['figure.figsize'] = (10, 6)
plt.rcParams['figure.dpi'] = 110

# The installed scikit-learn still supports SVC(probability=True) but warns; silence that specific notice
warnings.filterwarnings("ignore", message="The `probability` parameter was deprecated", category=FutureWarning)
np.random.seed(42)
print("Semi-supervised libraries imported successfully!")

# ---- Load the Breast Cancer Wisconsin dataset ----
data = load_breast_cancer()
X_full = data.data            # 30 continuous diagnostic features
y_full = data.target          # 0 = malignant, 1 = benign
feature_names = data.feature_names
target_names = data.target_names

# ---- Train / Test split (stratified to preserve class balance) ----
X_train_full, X_test, y_train_full, y_test = train_test_split(
    X_full, y_full, test_size=0.30, random_state=42, stratify=y_full
)

# ---- Standardize features (fit on train only to avoid test leakage) ----
scaler = StandardScaler()
X_train_full_scaled = scaler.fit_transform(X_train_full)
X_test_scaled = scaler.transform(X_test)

# ---- Simulate the semi-supervised setting: hide 85% of the training labels ----
unlabeled_ratio = 0.85
rng = np.random.RandomState(42)
random_unlabeled_points = rng.rand(len(y_train_full)) < unlabeled_ratio   # True -> mask

y_train_semi = np.copy(y_train_full)
y_train_semi[random_unlabeled_points] = -1   # -1 is the scikit-learn convention for "unlabeled"

n_labeled = np.sum(y_train_semi != -1)
n_unlabeled = np.sum(y_train_semi == -1)

print(f"Total Training Samples:   {len(y_train_full)}")
print(f" - Labeled Samples (15%): {n_labeled}")
print(f" - Unlabeled Samples (85%): {n_unlabeled}")
print(f"Holdout Test Samples:     {len(y_test)}")
\end{lstlisting}

\textbf{Listing 7.2 --- Experiment 1: Self-Training with an RBF-kernel
SVC}

\begin{lstlisting}[language=Python]
# ---- Experiment 1: Self-Training with an RBF-kernel SVC ----

# 1. Supervised baseline trained ONLY on the ~60 labeled samples
mask_labeled = (y_train_semi != -1)                 # boolean mask of the labeled subset
X_labeled = X_train_full_scaled[mask_labeled]
y_labeled = y_train_semi[mask_labeled]

base_svc = SVC(kernel='rbf', probability=True, C=1.0, random_state=42)
base_svc.fit(X_labeled, y_labeled)
y_pred_baseline_svc = base_svc.predict(X_test_scaled)
acc_baseline_svc = accuracy_score(y_test, y_pred_baseline_svc)
f1_baseline_svc = f1_score(y_test, y_pred_baseline_svc)

# 2. Semi-supervised self-training: base SVC + iterative pseudo-labeling above threshold 0.80
self_training_svc = SelfTrainingClassifier(
    estimator=SVC(kernel='rbf', probability=True, C=1.0, random_state=42),
    threshold=0.80,          # only accept pseudo-labels with confidence >= 0.80
    criterion='threshold',   # use the confidence-threshold selection rule
    max_iter=15,
    verbose=True             # print how many labels are added per iteration
)
self_training_svc.fit(X_train_full_scaled, y_train_semi)
y_pred_semi_svc = self_training_svc.predict(X_test_scaled)
acc_semi_svc = accuracy_score(y_test, y_pred_semi_svc)
f1_semi_svc = f1_score(y_test, y_pred_semi_svc)

# 3. Fully supervised "ceiling": train on 100% of the ground-truth labels
full_svc = SVC(kernel='rbf', probability=True, C=1.0, random_state=42)
full_svc.fit(X_train_full_scaled, y_train_full)
y_pred_full_svc = full_svc.predict(X_test_scaled)
acc_full_svc = accuracy_score(y_test, y_pred_full_svc)
f1_full_svc = f1_score(y_test, y_pred_full_svc)

# ---- Compare the three settings ----
print(f"\n--- SVC Evaluation Results ---")
print(f"Supervised Baseline (15% labeled): Accuracy = {acc_baseline_svc:.4f}, F1 = {f1_baseline_svc:.4f}")
print(f"Self-Training (Semi-supervised):   Accuracy = {acc_semi_svc:.4f}, F1 = {f1_semi_svc:.4f}")
print(f"Fully Supervised Ceiling (100%):   Accuracy = {acc_full_svc:.4f}, F1 = {f1_full_svc:.4f}")
\end{lstlisting}

\textbf{Listing 7.3 --- Inspect how self-training converged}

\begin{lstlisting}[language=Python]
# ---- Inspect how self-training converged ----
# n_iter_                 : number of self-training iterations performed
# termination_condition_  : why training stopped (e.g. 'no_change' = no new confident labels)
# transduction_           : final labels for every training sample (-1 = still unlabeled)
n_iter = self_training_svc.n_iter_
termination_condition = self_training_svc.termination_condition_
labeled_final_mask = (self_training_svc.transduction_ != -1)

print(f"Self-Training Convergence Summary:")
print(f" - Iterations completed: {n_iter}")
print(f" - Termination condition: {termination_condition}")
print(f" - Total samples labeled after self-training: {np.sum(labeled_final_mask)} / {len(y_train_full)}")
print(f" - Pseudo-labels newly added: {np.sum(labeled_final_mask) - n_labeled}")
\end{lstlisting}

\textbf{Listing 7.4 --- Compact comparison table (SVC)}

\begin{lstlisting}[language=Python]
# ---- Compact comparison table (SVC) ----
summary_data = [
    {"Model": "SVC (Baseline 15% Labeled)", "Accuracy": acc_baseline_svc, "F1-Score": f1_baseline_svc,
     "Training Samples": n_labeled},
    {"Model": "SVC (Self-Training Semi-Supervised)", "Accuracy": acc_semi_svc, "F1-Score": f1_semi_svc,
     "Training Samples": np.sum(self_training_svc.transduction_ != -1)},
    {"Model": "SVC (Fully Supervised Ceiling 100%)", "Accuracy": acc_full_svc, "F1-Score": f1_full_svc,
     "Training Samples": len(y_train_full)},
]
summary_df = pd.DataFrame(summary_data).set_index("Model")
display(summary_df.style.highlight_max(subset=['Accuracy', 'F1-Score'], color='lightgreen'))
\end{lstlisting}
